% ============================================================
%   LECTURE NOTES TEMPLATE
%   Supports: theorems, proofs, code, plots, geometry, figures
% ============================================================

\documentclass[11pt, a4paper]{article}

% ── Core packages ───────────────────────────────────────────
\usepackage[T1]{fontenc}
\usepackage[utf8]{inputenc}
\usepackage{lmodern}                   % Better fonts
\usepackage[margin=2.5cm]{geometry}    % Page margins
\usepackage{microtype}                 % Better text justification

% ── Math ────────────────────────────────────────────────────
\usepackage{amsmath, amssymb, amsthm, mathtools}
\usepackage{thmtools}                  % Extended theorem tools
\usepackage{bbm}                       % \mathbbm{1} indicator
\usepackage{bm}                        % Bold math \bm{}

% ── Colors ──────────────────────────────────────────────────
\usepackage[dvipsnames, svgnames, x11names]{xcolor}

% ── Define Color Palette ────────────────────────────────────
\definecolor{TheoremBG}    {RGB}{230, 245, 230}   % soft green
\definecolor{TheoremBorder}{RGB}{ 34, 139,  34}   % forest green
\definecolor{PropBG}       {RGB}{255, 224, 224}   % soft blue
\definecolor{PropBorder}   {RGB}{255, 123, 123}   % steel blue
\definecolor{CorolBG}      {RGB}{255, 245, 220}   % warm cream
\definecolor{CorolBorder}  {RGB}{205, 133,  63}   % peru
\definecolor{DefBG}        {RGB}{245, 230, 255}   % lavender
\definecolor{DefBorder}    {RGB}{128,   0, 128}   % purple
\definecolor{ExBG}         {RGB}{240, 248, 255}   % alice blue
\definecolor{ExBorder}     {RGB}{ 70, 130, 180}   % steel blue
\definecolor{RemBG}        {RGB}{255, 250, 240}   % floral white
\definecolor{RemBorder}    {RGB}{169, 169, 169}   % dark gray
\definecolor{LemmaBG}      {RGB}{230, 250, 250}   % light cyan
\definecolor{LemmaBorder}  {RGB}{ 32, 178, 170}   % light sea green
\definecolor{ProblemBG}    {RGB}{230, 240, 255}   % same as Prop
\definecolor{ProblemBorder}{RGB}{30, 90, 200}   % royal blue
\definecolor{CodeBG}       {RGB}{ 40,  44,  52}   % dark (One Dark)
\definecolor{CodeKey}      {RGB}{198, 120, 221}   % purple keyword
\definecolor{CodeStr}      {RGB}{152, 195, 121}   % green string
\definecolor{CodeCom}      {RGB}{ 92, 100, 112}   % gray comment
\definecolor{CodeNum}      {RGB}{209, 154, 102}   % orange number
\definecolor{LinkColor}    {RGB}{ 10,  80, 160}

% ── Theorem Environments (tcolorbox) ────────────────────────
\usepackage[most]{tcolorbox}
\tcbuselibrary{theorems, skins, breakable}

%  Macro for colored theorem boxes
\newcommand{\maketheorbox}[6]{%
  % #1 = env name, #2 = display name, #3 = BG color, #4 = border color
  % #5 = title text color, #6 = counter (share with: e.g. theorem)
  \tcbset{#1style/.style={%
    enhanced, breakable,
    colback=#3, colframe=#4,
    arc=4pt,
    fonttitle=\bfseries\color{white},
    attach boxed title to top left={yshift=-\tcboxedtitleheight/2, xshift=8pt},
    boxed title style={colback=#4, colframe=#4, sharp corners},
    sharp corners=south,
    boxrule=1pt,
    top=4mm, bottom=4mm,
  }}%
}

% ── Theorem / Proposition / etc. counters ───────────────────
\newcounter{thmcounter}[section]
\renewcommand{\thethmcounter}{\thesection.\arabic{thmcounter}}

% THEOREM
\newtcbtheorem[use counter=thmcounter]{theorem}{Theorem}{
  enhanced, breakable,
  colback=TheoremBG, colframe=TheoremBorder,
  arc=4pt,
  fonttitle=\bfseries\color{white},
  attach boxed title to top left={yshift=-\tcboxedtitleheight/2, xshift=8pt},
  boxed title style={colback=TheoremBorder, colframe=TheoremBorder,
    sharp corners, fontupper=\color{white}\bfseries},
  sharp corners=south, boxrule=1pt,
  top=5mm, bottom=4mm,
}{thm}

% PROPOSITION
\newtcbtheorem[use counter=thmcounter]{proposition}{Proposition}{
  enhanced, breakable,
  colback=PropBG, colframe=PropBorder,
  arc=4pt,
  fonttitle=\bfseries\color{white},
  attach boxed title to top left={yshift=-\tcboxedtitleheight/2, xshift=8pt},
  boxed title style={colback=PropBorder, colframe=PropBorder,
    sharp corners, fontupper=\color{white}\bfseries},
  sharp corners=south, boxrule=1pt,
  top=5mm, bottom=4mm,
}{prop}

% LEMMA
\newtcbtheorem[use counter=thmcounter]{lemma}{Lemma}{
  enhanced, breakable,
  colback=LemmaBG, colframe=LemmaBorder,
  arc=4pt,
  fonttitle=\bfseries\color{white},
  attach boxed title to top left={yshift=-\tcboxedtitleheight/2, xshift=8pt},
  boxed title style={colback=LemmaBorder, colframe=LemmaBorder,
    sharp corners, fontupper=\color{white}\bfseries},
  sharp corners=south, boxrule=1pt,
  top=5mm, bottom=4mm,
}{lem}

% COROLLARY
\newtcbtheorem[use counter=thmcounter]{corollary}{Corollary}{
  enhanced, breakable,
  colback=CorolBG, colframe=CorolBorder,
  arc=4pt,
  fonttitle=\bfseries\color{white},
  attach boxed title to top left={yshift=-\tcboxedtitleheight/2, xshift=8pt},
  boxed title style={colback=CorolBorder, colframe=CorolBorder,
    sharp corners, fontupper=\color{white}\bfseries},
  sharp corners=south, boxrule=1pt,
  top=5mm, bottom=4mm,
}{cor}

% DEFINITION
\newtcbtheorem[use counter=thmcounter]{definition}{Definition}{
  enhanced, breakable,
  colback=DefBG, colframe=DefBorder,
  arc=4pt,
  fonttitle=\bfseries\color{white},
  attach boxed title to top left={yshift=-\tcboxedtitleheight/2, xshift=8pt},
  boxed title style={colback=DefBorder, colframe=DefBorder,
    sharp corners, fontupper=\color{white}\bfseries},
  sharp corners=south, boxrule=1pt,
  top=5mm, bottom=4mm,
}{def}

% EXAMPLE
\newtcbtheorem[use counter=thmcounter]{example}{Example}{
  enhanced, breakable,
  colback=ExBG, colframe=ExBorder,
  arc=4pt,
  fonttitle=\bfseries\color{white},
  attach boxed title to top left={yshift=-\tcboxedtitleheight/2, xshift=8pt},
  boxed title style={colback=ExBorder, colframe=ExBorder,
    sharp corners, fontupper=\color{white}\bfseries},
  sharp corners=south, boxrule=0.8pt,
  top=5mm, bottom=4mm,
}{ex}

% REMARK  (lighter, no box title)
\newtcbtheorem[use counter=thmcounter]{remark}{Remark}{
  enhanced, breakable,
  colback=RemBG, colframe=RemBorder,
  arc=4pt,
  fonttitle=\bfseries\color{white},
  attach boxed title to top left={yshift=-\tcboxedtitleheight/2, xshift=8pt},
  boxed title style={colback=RemBorder, colframe=RemBorder,
    sharp corners, fontupper=\color{white}\bfseries},
  sharp corners=south, boxrule=0.6pt,
  top=5mm, bottom=4mm,
}{rem}

% PROBLEM
\newtcbtheorem[use counter=thmcounter]{problem}{Problem}{
  enhanced, breakable,
  colback=ProblemBG, colframe=ProblemBorder,
  arc=4pt,
  fonttitle=\bfseries\color{white},
  attach boxed title to top left={yshift=-\tcboxedtitleheight/2, xshift=8pt},
  boxed title style={colback=ProblemBorder, colframe=ProblemBorder,
    sharp corners, fontupper=\color{white}\bfseries},
  sharp corners=south, boxrule=1pt,
  top=5mm, bottom=4mm,
}{prob}

% ── PROOF style (white background, left gray bar, italic) ───
% Matches the screenshot: "Proof of Theorem X.X" header, plain body
\tcolorboxenvironment{proof}{
  blanker, breakable,
  left=4mm,
  before skip=6pt, after skip=10pt,
  borderline west={2pt}{0pt}{gray!60},
}
% Override proof header style
\renewcommand{\proofname}{\normalfont\color{gray}\small\scshape Proof}

% ── Code listings (listings + tcolorbox) ────────────────────
\usepackage{listings}

\lstdefinestyle{CppStyle}{
  language=C++,
  backgroundcolor=\color{CodeBG},
  basicstyle=\ttfamily\small\color{white},
  keywordstyle=\color{CodeKey}\bfseries,
  stringstyle=\color{CodeStr},
  commentstyle=\color{CodeCom}\itshape,
  numberstyle=\tiny\color{CodeNum},
  numbers=left, numbersep=8pt,
  stepnumber=1,
  breaklines=true, breakatwhitespace=false,
  tabsize=4,
  showstringspaces=false,
  morekeywords={auto, nullptr, constexpr, override, final,
                uint8_t, uint16_t, uint32_t, uint64_t,
                int8_t, int16_t, int32_t, int64_t},
  frame=none,
}

\lstdefinestyle{PythonStyle}{
  language=Python,
  backgroundcolor=\color{CodeBG},
  basicstyle=\ttfamily\small\color{white},
  keywordstyle=\color{CodeKey}\bfseries,
  stringstyle=\color{CodeStr},
  commentstyle=\color{CodeCom}\itshape,
  numberstyle=\tiny\color{CodeNum},
  numbers=left, numbersep=8pt,
  stepnumber=1,
  breaklines=true, breakatwhitespace=false,
  tabsize=4,
  showstringspaces=false,
  morekeywords={self, None, True, False, as, with, yield,
                lambda, import, from, pass, assert, raise,
                except, finally, elif},
  frame=none,
}

% Wrap lstlisting in a tcolorbox for rounded corners + title
\newtcblisting{cppcode}[1][]{
  listing only, listing style=CppStyle,
  colback=CodeBG, colframe=CodeBG!80!black,
  fonttitle=\ttfamily\small\color{white},
  title=\faCode\ C++,            % remove \faCode if not using fontawesome
  sharp corners, boxrule=0.5pt,
  #1
}

\newtcblisting{pycode}[1][]{
  listing only, listing style=PythonStyle,
  colback=CodeBG, colframe=CodeBG!80!black,
  fonttitle=\ttfamily\small\color{white},
  title=Python,
  sharp corners, boxrule=0.5pt,
  #1
}

% Fallback plain code block (language-agnostic)
\newtcblisting{code}[2][]{
  listing only,
  listing options={
    style=PythonStyle,
    language=#2,
  },
  colback=CodeBG, colframe=CodeBG!80!black,
  fonttitle=\ttfamily\small\color{white},
  title=#2,
  sharp corners, boxrule=0.5pt,
  #1
}

% ── Plots and Geometry (pgfplots + tikz) ────────────────────
\usepackage{tikz}
\usepackage{pgfplots}
\pgfplotsset{compat=1.18}
\usetikzlibrary{
  arrows.meta,         % Better arrow tips
  calc,                % Coordinate calculations
  shapes.geometric,    % circles, rectangles, etc.
  positioning,         % node placement
  decorations.markings,% marks on paths
  intersections,       % path intersections
  patterns,            % hatching
  angles, quotes,      % angle labeling
  3d,                  % 3D axes
}

% Convenience: TikZ picture wrapper with caption
\usepackage{float}
\usepackage{caption}
\usepackage{subcaption}

% ── Figures / Images ────────────────────────────────────────
\usepackage{graphicx}
\usepackage{wrapfig}   % Text-wrapped figures

% Figure box style (tcolorbox)
\newtcolorbox{figurebox}[2][]{
  enhanced, breakable,
  colback=white, colframe=gray!50,
  title=#2,
  fonttitle=\bfseries\small,
  boxed title style={colback=gray!20, colframe=gray!50},
  attach boxed title to top left={yshift=-\tcboxedtitleheight/2, xshift=8pt},
  sharp corners, boxrule=0.8pt,
  #1
}

% ── Tables ──────────────────────────────────────────────────
\usepackage{booktabs}   % \toprule, \midrule, \bottomrule
\usepackage{array}
\usepackage{multirow}
\usepackage{longtable}  % Tables spanning multiple pages

% ── Hyperlinks & PDF metadata ───────────────────────────────
\usepackage{hyperref}
\hypersetup{
  colorlinks = true,
  linkcolor  = LinkColor,
  citecolor  = LinkColor,
  urlcolor   = LinkColor,
  pdftitle   = {Lecture Notes},
  pdfauthor  = {Your Name},
}

% ── Table of Contents style ─────────────────────────────────
\usepackage{titletoc}
\usepackage{titlesec}

% Section numbering depth
\setcounter{tocdepth}{2}
\setcounter{secnumdepth}{3}

% Section title styling
\titleformat{\section}
  {\large\bfseries}
  {\thesection}{1em}{}[\titlerule]

\titleformat{\subsection}
  {\normalsize\bfseries}
  {\thesubsection}{1em}{}

% ── Bibliography ────────────────────────────────────────────
\usepackage[backend=biber, style=alphabetic]{biblatex}
% \addbibresource{references.bib}   % uncomment and add your .bib

% ── Miscellaneous Useful Packages ───────────────────────────
\usepackage{enumitem}          % Customise lists
\usepackage{todonotes}         % \todo{} margin notes
\usepackage{cleveref}          % \cref{} smart cross-refs
\usepackage{siunitx}           % SI units: \SI{9.8}{\metre\per\second}
\usepackage{physics}           % \dd, \grad, \curl, \laplacian …
\usepackage{mathrsfs}          % \mathscr{F} script font

% ── Convenience Math Macros ─────────────────────────────────
\newcommand{\R}{\mathbb{R}}
\newcommand{\N}{\mathbb{N}}
\newcommand{\Z}{\mathbb{Z}}
\newcommand{\Q}{\mathbb{Q}}
\newcommand{\C}{\mathbb{C}}
\newcommand{\E}{\mathbb{E}}
\newcommand{\Prob}{\mathbb{P}}
\newcommand{\one}{\mathbbm{1}}          % indicator
\newcommand{\eps}{\varepsilon}
% \newcommand{\norm}[1]{\left\lVert#1\right\rVert}
% \newcommand{\abs}[1]{\left\lvert#1\right\rvert}
\newcommand{\inner}[2]{\left\langle#1,\,#2\right\rangle}
\newcommand{\set}[1]{\left\{#1\right\}}
\newcommand{\floor}[1]{\left\lfloor#1\right\rfloor}
\newcommand{\ceil}[1]{\left\lceil#1\right\rceil}
\DeclareMathOperator*{\argmin}{arg\,min}
\DeclareMathOperator*{\argmax}{arg\,max}
\DeclareMathOperator{\sgn}{sgn}
% \DeclareMathOperator{\tr}{tr}
% \DeclareMathOperator{\rank}{rank}
\DeclareMathOperator{\diag}{diag}
\DeclareMathOperator{\supp}{supp}

% ── Header / Footer ─────────────────────────────────────────
\usepackage{fancyhdr}
\pagestyle{fancy}
\fancyhf{}
\fancyhead[L]{\small\textit{\CourseName}}
\fancyfoot[C]{\thepage}
\renewcommand{\headrulewidth}{0.4pt}

% ── Metadata Commands (fill these per lecture) ──────────────
\newcommand{\CourseName}{Topology \& Multivariable Calculus}
\newcommand{\CourseCode}{FMA\_2F002\_EP}
\newcommand{\University}{Ecole Polytechnique}
\newcommand{\Lecturer}{Prof. Luc Pirio}
\newcommand{\StudentName}{Nikolay Kharovskiy}
\newcommand{\Semester}{Fall 2026}

% ============================================================
%   DOCUMENT BEGIN
% ============================================================
\begin{document}

% ── Title Page ──────────────────────────────────────────────
\begin{titlepage}
  \centering
  \vspace*{2cm}

  {\Huge\bfseries \CourseName \par}
  \vspace{0.5cm}
  {\large \CourseCode \par}

  \vspace{1.5cm}
  \hrule height 1pt
  \vspace{0.4cm}
  {\Large\bfseries Lecture Notes \par}
  \vspace{0.4cm}
  \hrule height 1pt

  \vspace{2cm}

  \begin{tabular}{rl}
    \textbf{University:} & \University \\[4pt]
    \textbf{Lecturer:}   & \Lecturer   \\[4pt]
    \textbf{Student:}    & \StudentName \\[4pt]
    \textbf{Semester:}   & \Semester   \\[4pt]
  \end{tabular}

  \vfill

  {\small These are personal lecture notes and may contain errors.
         For the official course material consult the lecturer.}
\end{titlepage}

% ── Table of Contents ───────────────────────────────────────
\tableofcontents
\newpage

% ── Lecture 1: Norms in vector spaces & Sequences ───────────
\section{Norms in Vector Spaces}
\subsection{Setting}
Let $E, F, G, \dots$ be vector spaces over $\mathbb{K} = \mathbb{R}$ or $\mathbb{C}$.
\begin{itemize}
    \item If $\mathbb{K} = \mathbb{R}$, for $x \in \mathbb{R}$, $|x| = \max(x, -x)$.
    \item If $\mathbb{K} = \mathbb{C}$, for $z = x + iy$, we have $|z| = \sqrt{x^2 + y^2} = l$. In polar form, $z = l e^{i\theta}$ where $l \ge 0$ and $\theta \in [0, 2\pi)$.
\end{itemize}

\subsection{Norms}
\begin{definition}{Norm}{norm_def}
A norm on $E$ is a map $N: E \to \mathbb{R}_{\ge 0} = \mathbb{R}^+ = [0, +\infty)$ satisfying the following 3 properties:
\begin{enumerate}
    \item \textbf{Triangular inequality:} $\forall x,y \in E$, $N(x+y) \le N(x) + N(y)$.
    \item \textbf{Absolute homogeneity:} $\forall \lambda \in \mathbb{K}, \forall x \in E$, $N(\lambda x) = |\lambda| N(x)$.
    \item \textbf{Positive definiteness:} $N(x) = 0 \Leftrightarrow x = 0$.
\end{enumerate}
\end{definition}

\begin{remark}{}{}
Positive definiteness is equivalent to the fact that $N$ "separates points". That is, $\forall x, y \in E$, $N(x-y) = 0$ if and only if $x = y$.
\end{remark}

\begin{proposition}{Reverse Triangular Inequality}{rev_tri}
The triangular inequality is equivalent to the reverse triangular inequality:
$\forall x, y \in E$, $\big|N(x) - N(y)\big| \le N(x-y)$.
\end{proposition}

\begin{proof}
Let $x, y \in E$. Then $N(x) = N(x - y + y) \le N(x-y) + N(y)$, which implies $N(x) - N(y) \le N(x-y)$.
Similarly, $N(y) - N(x) \le N(y-x)$.
Since $N(x-y) = |-1|N(y-x)$, we obtain $\big|N(x) - N(y)\big| \le N(x-y)$.
\end{proof}

\begin{example}{Common Norms}{common_norms}
\begin{itemize}
    \item $|\cdot|$ is a norm on $\mathbb{R}$.
    \item On $\displaystyle E = \mathbb{R}^d$, Sup Norm $\displaystyle ||x||_\infty = \max_{i=1}^d |x_i|$.
    \item Euclidean norm on $\displaystyle \mathbb{R}^d$: $\displaystyle||x||_2 = \sqrt{\sum_{i=1}^d x_i^2}$.
    \item On $F = C^0([0,1], \mathbb{R})$, Sup Norm $\displaystyle ||f||_\infty = \sup_{x \in [0,1]} |f(x)|$.
    \item On $F$, $L_1$ norm $\displaystyle||f||_{L^1} = \int_0^1 |f(x)| dx$.
    \item On $F$, $L_2$ norm $\displaystyle||f||_{L^2} = \sqrt{\int_0^1 |f(x)|^2 dx}$.
\end{itemize}
\end{example}

\begin{lemma}{Norms on $\mathbb{R}$}{norm_r}
Let $N: \mathbb{R} \to \mathbb{R}_{\ge 0}$ be a norm. Then $\exists \alpha > 0$ such that $N(\cdot) = \alpha |\cdot|$.
\end{lemma}
\begin{proof}
Let $x \in \mathbb{R}$. Then $N(x) = N(x \cdot 1) = |x|N(1)$. Let $\alpha = N(1) > 0$.
\end{proof}

\begin{lemma}{Euclidean Norm on $\mathbb{R}^d$}{norm_euclidean}
$||\cdot||_2$ is a norm on $E = \mathbb{R}^d$.
\end{lemma}
\begin{proof}
Absolute homogeneity is obvious.
For positive definiteness: $||x||_2 = 0 \Rightarrow \sum_{i=1}^d x_i^2 = 0 \Rightarrow x_i = 0$ for all $i = 1, \dots, d \Rightarrow x = 0_E$.
For the triangular inequality, we use the Euclidean scalar product $\langle x, y \rangle = \sum_{i=1}^d x_i y_i$. This product is symmetric, $\mathbb{R}$-bilinear, and $\langle x, x \rangle = 0 \Rightarrow x = 0_E$.
We have $||x+y||_2 = \sqrt{\langle x+y, x+y \rangle} = \sqrt{||x||_2^2 + ||y||_2^2 + 2\langle x, y \rangle}$.
By the Cauchy-Schwarz inequality, $|\langle x, y \rangle| \le ||x||_2 ||y||_2$. Thus, $||x+y||_2 \le \sqrt{||x||_2^2 + 2||x||_2 ||y||_2 + ||y||_2^2} = ||x||_2 + ||y||_2$.
\end{proof}

\begin{lemma}{Sup Norm on $C^0([0,1], \mathbb{R})$}{norm_sup}
$||\cdot||_\infty$ is a norm on $F = C^0([0,1], \mathbb{R})$.
\end{lemma}
\begin{proof}
For any $f \in F$, $f$ is bounded.
Positive homogeneity: $||\lambda f||_\infty = \sup |\lambda f(x)| = |\lambda| ||f||_\infty$.
Positive definiteness: If $||f||_\infty = 0$, then $0 \le |f(x)| \le \sup |f(x)| = 0$, which implies $f(x) = 0$ for all $x$, so $f = 0_F$.
Triangular inequality: $|(f+g)(x)| = |f(x)+g(x)| \le |f(x)| + |g(x)| \le ||f||_\infty + ||g||_\infty$. Taking the supremum yields the result.
\end{proof}

\subsection{Distances}
\begin{definition}{Distance}{dist_def}
A distance on a set $X \neq \emptyset$ is a map $d: X^2 \to \mathbb{R}^+$ satisfying:
\begin{enumerate}
    \item \textbf{Triangular inequality:} $\forall x,y,z \in X$, $d(x, z) \le d(x, y) + d(y, z)$.
    \item \textbf{Symmetry:} $\forall x,y \in X$, $d(x, y) = d(y, x)$.
    \item \textbf{Positive definiteness:} $\forall x,y \in X$, $d(x, y) = 0 \Leftrightarrow x = y$.
\end{enumerate}
\end{definition}

\begin{example}{Discrete Distance}{discrete_dist}
Let $X \neq \emptyset$. The discrete distance is defined as $d(x,y) = 0$ if $x = y$, and $d(x,y) = 1$ if $x \neq y$.
\end{example}

\begin{proposition}{Distance Induced by a Norm}{dist_norm}
Let $(E, N)$ be a normed vector space (NVS) over $\mathbb{K}$. The map $d_N: E^2 \to \mathbb{R}^+$ given by $d(x,y) = N(x-y)$ is a distance on $E$.
\end{proposition}
\begin{proof}
Symmetry: $d(x,y) = N(x-y) = N(y-x) = d(y,x)$.
Positive definiteness: $d(x,y) = 0 \Leftrightarrow N(x-y) = 0 \Rightarrow x-y = 0 \Rightarrow x = y$.
Triangular inequality: $d(x,z) = N(x-z) = N(x-y+y-z) \le N(x-y) + N(y-z) = d(x,y) + d(y,z)$.
\end{proof}

\subsection{Balls and Convexity}
\begin{definition}{Balls}{ball_def}
Let $(X, d)$ be a metric space. For $x_0 \in X$ and $r \in \mathbb{R}^+$, we define:
\begin{itemize}
    \item \textbf{Open ball:} $B(x_0, r) = \{x \in X \mid d(x_0, x) < r\}$.
    \item \textbf{Closed ball:} $\overline{B}(x_0, r) = \{x \in X \mid d(x_0, x) \le r\}$.
\end{itemize}
In $E$, $B(0, 1)$ is defined as the unit ball.
\end{definition}

\begin{definition}{Segments and Convexity}{convex_def}
For $x, y \in E$, the closed segment is $[x,y] = \{tx + (1-t)y \mid t \in [0, 1]\}$. The open segment is $(x,y) = [x,y] \setminus \{x, y\}$.
A subset $C \subseteq E$ is \textbf{convex} if for all $x, y \in C$, the segment $[x,y] \subseteq C$.
\end{definition}

\begin{proposition}{Convex Subsets}{convex_prop}
\begin{enumerate}
    \item Convex subsets of $\mathbb{R}$ are intervals.
    \item Balls in a NVS are convex.
\end{enumerate}
\end{proposition}
\begin{proof}
We prove (2) for $B = B(x_0, r)$ with $r > 0$.
Let $x, y \in B$, meaning $N(x-x_0) < r$ and $N(y-x_0) < r$.
Let $z = tx + (1-t)y \in [x,y]$ with $t \in [0,1]$.
We have $d(z, x_0) = N(z-x_0) = N(t(x-x_0) + (1-t)(y-x_0))$.
By the triangular inequality, $d(z, x_0) \le tN(x-x_0) + (1-t)N(y-x_0) < tr + (1-t)r = r$.
Thus, $z \in B$, which implies $[x,y] \subseteq B$, meaning $B$ is convex.
\end{proof}

\section{Sequences}
\begin{definition}{Convergence of Sequences}{seq_conv}
Let $\underline{x} = (x_n)_{n \in \mathbb{N}}$ be a sequence in a $\mathbb{K}$-NVS $E$.
We say $x_n \to x_\infty$ in $E$ (wrt $N$) if $N(x_n - x_\infty) \to 0$ in $\mathbb{R}$.
We denote this as $x_n \to x_\infty$ or $\lim x_n = x_\infty$. If a sequence does not converge, it diverges.
\end{definition}

\begin{example}{Convergence in Different Norms}{seq_ex}
Consider $E = C^0([0,1], \mathbb{R})$ and the sequence of functions $f_n(x) = x^n$.
\begin{itemize}
    \item For the sup norm $||\cdot||_\infty$, we have $||f_n||_\infty = 1$ for all $n$, so it does not converge to $0_E$.
    \item For the $L^1$ norm, we have $||f_n||_{L^1} = \int_0^1 x^n dx = \frac{1}{n+1} \to 0$. Thus $f_n \to 0_E$ for $||\cdot||_{L^1}$.
\end{itemize}
\end{example}

\begin{proposition}{Uniqueness of the Limit}{limit_unique}
The limit of a convergent sequence is unique.
\end{proposition}
\begin{proof}
Let $(x_n)_{n \in \N}$ be a sequence in $E$ converging to both $x_\infty \in E$ and $x_\infty' \in E$.
By the triangular inequality, $0 \le N(x_\infty - x_\infty') \le N(x_\infty - x_n) + N(x_n - x_\infty')$.
Since $x_n \to x_\infty$ and $x_n \to x_\infty'$, both $N(x_\infty - x_n) \to 0$ and $N(x_n - x_\infty') \to 0$ as $n \to \infty$. By the squeeze theorem, $N(x_\infty - x_\infty') = 0$, which implies $x_\infty = x_\infty'$.
\end{proof}

\begin{proposition}{Continuity of the Norm}{norm_cont}
\begin{enumerate}
    \item If $x_n \to x_\infty$ in $E$, then $N(x_n) \to N(x_\infty)$ in $\R$.
    \item Hence, $(N(x_n))_{n \in \N}$ is a bounded sequence, and $N(x_\infty) \le \sup_{n \in \N} N(x_n)$.
\end{enumerate}
\end{proposition}
\begin{proof}
By the reverse triangular inequality, $0 \le |N(x_n) - N(x_\infty)| \le N(x_n - x_\infty)$.
Since $x_n \to x_\infty$, we have $N(x_n - x_\infty) \to 0$. By the squeeze theorem, $|N(x_n) - N(x_\infty)| \to 0$, meaning $N(x_n) \to N(x_\infty)$.
The second part follows from the fact that any convergent sequence $(u_n)_{n \in \N}$ in $\R$ is bounded (i.e., $\sup_{n \in \N} |u_n| < +\infty$). Its limit satisfies $|u_\infty| \le \sup_{n \in \N} |u_n|$.
\end{proof}

\begin{proposition}{Algebraic Operations on Limits}{limit_algebra}
Let $(x_n)_{n \in \N}$ and $(y_n)_{n \in \N}$ be sequences in $E$, and let $(\lambda_n)_{n \in \N}$ be a sequence in $\mathbb{K}$. Assume:
\begin{enumerate}
    \item $x_n \to x_\infty$ in $E$
    \item $y_n \to y_\infty$ in $E$
    \item $\lambda_n \to \lambda_\infty$ in $\mathbb{K}$
\end{enumerate}
Then, $x_n + y_n \to x_\infty + y_\infty$ in $E$, and $\lambda_n y_n \to \lambda_\infty y_\infty$ in $E$.
\end{proposition}
\begin{proof}
For $n \in \N$, set $z_n = x_n + \lambda_n y_n$ and $z_\infty = x_\infty + \lambda_\infty y_\infty$. We want to show $N(z_n - z_\infty) \to 0$.
By the triangular inequality, $0 \le N(z_n - z_\infty) \le N(x_n - x_\infty) + N(\lambda_n y_n - \lambda_\infty y_\infty)$.
By adding and subtracting $\lambda_\infty y_n$, we get $N(\lambda_n y_n - \lambda_\infty y_\infty) \le |\lambda_n - \lambda_\infty| N(y_n) + |\lambda_\infty| N(y_n - y_\infty)$.
Since $y_n \to y_\infty$, the sequence $(N(y_n))_{n \in \N}$ is bounded by some $M > 0$.
Thus, $N(\lambda_n y_n - \lambda_\infty y_\infty) \le M |\lambda_n - \lambda_\infty| + |\lambda_\infty| N(y_n - y_\infty)$.
As $n \to \infty$, $|\lambda_n - \lambda_\infty| \to 0$ and $N(y_n - y_\infty) \to 0$, so $N(\lambda_n y_n - \lambda_\infty y_\infty) \to 0$.
Consequently, $N(z_n - z_\infty) \to 0$, which completes the proof.
\end{proof}

\begin{definition}{Subsequence and Adherence Value}{subseq_def}
\begin{itemize}
    \item A sequence $(y_n)_{n \in \N}$ is a \textbf{subsequence} (or extraction) of $(x_n)_{n \in \N}$ if there exists a strictly increasing function $\varphi: \N \to \N$ such that $y_n = x_{\varphi(n)}$.
    \item An element $l \in E$ is an \textbf{adherence value} (or accumulation point) of $(x_n)_{n \in \N}$ if there exists a strictly increasing extraction $\varphi: \N \to \N$ such that $x_{\varphi(n)} \to l$.
\end{itemize}
\end{definition}

\begin{problem}{Accumulation Points of $\Q$}{q_accum}
The set of rational numbers $\Q$ is countable, so it can be written as a sequence $\Q = \{q_0, q_1, \dots\}$. What are the accumulation points of the sequence $(q_n)_{n \in \N}$?
\end{problem}
\begin{remark}{}{}
Because $\Q$ is dense in $\R$, every real number is an accumulation point. Thus, the set of accumulation points is all of $\R$.
\end{remark}

% ── Lecture 2 ───────────────────────────────────────────────
\section{Lecture 2: Open \& Closed Subsets in Normed Vector Spaces}
Let $(E, \Vert \cdot \Vert)$ be a fixed NVS over $\mathbb{K}$.

\subsection{Open Subsets}
\begin{definition}{Open Subset}{open_def}
A subset $U \subset E$ is \textbf{open} if for all $u \in U$, there exists $\eps > 0$ such that $B(u, \eps) \subset U$.
\end{definition}

\begin{example}{Open Sets}{open_ex}
\begin{itemize}
    \item $\emptyset$ and $E$ are open subsets.
    \item In $\R$, open intervals $(a, b)$ are open subsets.
\end{itemize}
\end{example}

\begin{proposition}{Open Balls are Open Sets}{open_ball_open}
In $(E, \Vert \cdot \Vert)$, open balls are open subsets.
\end{proposition}
\begin{proof}
Let $B = B(x_0, r) \subset E$ with $r > 0$. Let $y \in B$.
We have $d(x_0, y) = \Vert x_0 - y \Vert = \delta < r$.
Choose $\eps_y = r - \delta = r - \Vert x_0 - y \Vert > 0$. We claim $B(y, \eps_y) \subset B$.
Let $z \in B(y, \eps_y)$. By the triangular inequality:
$d(x_0, z) \le d(x_0, y) + d(y, z) = \Vert x_0 - y \Vert + \Vert y - z \Vert < \delta + \eps_y = r$.
Therefore, $z \in B$, which implies $B(y, \eps_y) \subset B$. Thus, $B$ is open.
\end{proof}

\begin{proposition}{Union of Open Balls}{open_union_balls}
The open subsets of $E$ are exactly the unions of open balls.
\end{proposition}
\begin{proof}
Let $U \subset E$ be an open subset.
By definition, for every $u \in U$, there exists $\eps_u > 0$ such that $B(u, \eps_u) \subset U$.
Then $U \subset \bigcup_{u \in U} B(u, \eps_u) \subset U$.
Thus, $U = \bigcup_{u \in U} B(u, \eps_u)$, which is a union of open balls.
\end{proof}

\subsection{Closed Subsets}
\begin{definition}{Closed Subset}{closed_def}
A subset $F \subset E$ is \textbf{closed} if its complement $F^c = E \setminus F$ is open.
\end{definition}

\begin{example}{Closed Sets}{closed_ex}
$\emptyset$ and $E$ are closed subsets of $(E, \Vert \cdot \Vert)$.
\end{example}

\begin{proposition}{Properties of Open and Closed Sets}{open_closed_prop}
\begin{enumerate}
    \item Arbitrary unions of open subsets are open.
    \item Finite intersections of open subsets are open.
    \item Arbitrary intersections of closed subsets are closed.
    \item Finite unions of closed subsets are closed.
\end{enumerate}
\end{proposition}
\begin{remark}{}{}
An infinite intersection of open sets is not necessarily open. For example, in $\R$, the intersection $\bigcap_{n \ge 1} (-\frac{1}{n}, \frac{1}{n}) = \{0\}$ is not open.
\end{remark}
\begin{proof}
Let us prove the properties for open sets:
\begin{enumerate}
    \item Let $(O_i)_{i \in I}$ be a family of open subsets. Let $\mathcal{O} = \bigcup_{i \in I} O_i$. If $x \in \mathcal{O}$, then there exists $i_0 \in I$ such that $x \in O_{i_0}$. Since $O_{i_0}$ is open, there exists $\eps > 0$ such that $B(x, \eps) \subset O_{i_0} \subset \mathcal{O}$. Thus, $\mathcal{O}$ is open.
    \item Assume $I$ is finite and non-empty. Let $\mathcal{O} = \bigcap_{i \in I} O_i$. Let $u \in \mathcal{O}$. Then for all $i \in I$, $u \in O_i$. Since each $O_i$ is open, there exists $\eps_i > 0$ such that $B(u, \eps_i) \subset O_i$. Because $I$ is finite, $\eps = \min_{i \in I} \eps_i > 0$. Then $B(u, \eps) \subset \bigcap_{i \in I} B(u, \eps_i) \subset \mathcal{O}$. Thus, $\mathcal{O}$ is open.
\end{enumerate}
The properties for closed sets follow symmetrically.
\end{proof}

\subsection{Sequential Characterizations}
\begin{theorem}{Sequential Characterization of Closed Subsets}{seq_closed}
For $F \subset E$, the following assertions are equivalent:
\begin{enumerate}
    \item $F$ is closed.
    \item For every sequence $(x_n)_{n \in \N}$ in $F$ converging in $E$, its limit $\lim x_n$ belongs to $F$.
\end{enumerate}
\end{theorem}
\begin{proof}
\textbf{($1 \Rightarrow 2$):} Assume $F$ is closed. We proceed by contradiction. Assume there exists a sequence $(x_n)_{n \in \N}$ in $F$ such that $x_n \to x_\infty \in E$, but $x_\infty \notin F$. Then $x_\infty \in F^c$.
Since $F$ is closed, $F^c$ is open. Therefore, there exists $\eps > 0$ such that $B(x_\infty, \eps) \subset F^c$.
Because $x_n \to x_\infty$, there exists $N_\eps$ such that for all $n \ge N_\eps$, $x_n \in B(x_\infty, \eps) \subset F^c$.
This implies $x_n \in F \cap F^c = \emptyset$, which is a contradiction. Thus, $x_\infty \in F$.

\textbf{($\neg 1 \Rightarrow \neg 2$):} Assume $F$ is not closed. Then $F^c$ is not open.
This means there exists $u^* \in F^c$ such that for all $\eps > 0$, $B(u^*, \eps) \not\subset F^c$, which is equivalent to saying $B(u^*, \eps) \cap F \neq \emptyset$.
For each $n \ge 1$, choose $x_n \in F \cap B(u^*, \frac{1}{n})$. This constructs a sequence $(x_n)_{n \in \N}$ in $F$.
Since $d(x_n, u^*) < \frac{1}{n}$, we have $x_n \to u^*$ as $n \to \infty$. However, $u^* \notin F$, negating the second assertion.
\end{proof}

\begin{theorem}{Sequential Characterization of Open Subsets}{seq_open}
For $U \subset E$, the following are equivalent:
\begin{enumerate}
    \item $U$ is open.
    \item For every $u \in U$ and for every sequence $(x_n)_{n \in \N}$ converging to $u$ in $E$, there exists $N \in \N$ such that for all $n \ge N$, $x_n \in U$.
\end{enumerate}
\end{theorem}

\subsection{Connectedness and Interior/Closure}
\begin{proposition}{Clopen Subsets}{clopen}
If $X \subset E$ (or $\R$) is both closed and open, then $X = \emptyset$ or $X = E$.
\end{proposition}
\begin{proof}
Assume for contradiction that $\emptyset \subsetneq X \subsetneq E$.
Then there exists $x \in X$ and $y \in X^c$.
Define $I = \{t \in [0, 1] \mid x + t(y - x) \in X\}$. Since $x \in X$, $0 \in I$. Since $y \notin X$, $1 \notin I$.
Let $z = \sup I$. Clearly, $z \le 1$.
By definition of supremum, there exists a sequence $(t_n)_{n \in \N}$ in $I$ such that $t_n \to z$.
The sequence $x_{t_n} = x + t_n(y - x)$ lies in $X$. Since $X$ is closed, the limit $x_z = x + z(y - x)$ must belong to $X$. Thus, $z \in I$, which implies $z < 1$ (since $1 \notin I$).
Because $X$ is open and $x_z \in X$, there exists $\eps > 0$ such that $B(x_z, \eps) \subset X$.
This means for sufficiently small $\delta > 0$, $x_{z+\delta} \in X$, implying $z+\delta \in I$.
This contradicts $z = \sup I$. Therefore, $X = \emptyset$ or $X = E$.
\end{proof}

\begin{definition}{Interior and Closure}{int_closure}
Let $X \subset E$.
\begin{itemize}
    \item An element $x \in X$ is an \textbf{interior point} of $X$ if there exists $\eps > 0$ such that $B(x, \eps) \subset X$.
    \item The \textbf{interior} of $X$, denoted $\mathring{X}$, is the set of all interior points of $X$. It is the largest open subset contained in $X$, given by $\displaystyle \mathring{X} = \bigcup_{\substack{O \subset X \\ O \text{ open}}} O$.
    \item An element $x \in E$ is \textbf{adherent} to $X$ if for every $\eps > 0$, $B(x, \eps) \cap X \neq \emptyset$.
    \item The \textbf{topological closure} of $X$, denoted $\overline{X}$, is the set of all adherent points to $X$. It is the smallest closed subset containing $X$, given by $\displaystyle \overline{X} = \bigcap_{\substack{F \supset X \\ F \text{ closed}}} F$.
\end{itemize}
\end{definition}

\begin{proposition}{Closure of a Ball}{ball_closure}
For $x_* \in E$ and $r > 0$, we have $\overline{B(x_*, r)} = \overline{B}(x_*, r)$, where $\overline{B(x_*, r)}$ is the topological closure of the open ball $B(x_*, r)$ and $\overline{B}(x_*, r)$ is the closed ball.
\end{proposition}
\begin{remark}{}{}
This is false in a general metric space.
For example, let $X$ be a set with at least two elements, and let $d$ be the discrete metric ($d(x,y) = 1, \forall x \neq y$).
Then for all $x \in X$, $\overline{B(x, 1)} = \{x\} \subsetneq X$.
\end{remark}
\begin{proof}
First, we show $\overline{B(x_*, r)} \subset \overline{B}(x_*, r)$.
The closed ball $\overline{B}(x_*, r)$ is a closed set, because its complement $U = E \setminus \overline{B}(x_*, r)$ is open.
To see this, let $y \in U$, so $\delta = d(x_*, y) > r$.
Let $\eps \in (0, \delta - r)$. By the reverse triangular inequality, $B(y, \eps) \subset U$, so $U$ is open.
Since $\overline{B}(x_*, r)$ is closed and contains $B(x_*, r)$, it must contain its closure, so $\overline{B(x_*, r)} \subset \overline{B}(x_*, r)$.

Now we prove the opposite inclusion. Let $x \in \overline{B}(x_*, r)$.
If $d(x_*, x) = N(x_* - x) < r$, then $x \in B(x_*, r) \subset \overline{B(x_*, r)}$.
Otherwise, $d(x_*, x) = r$.
We use the sequential characterization of the closure: $x \in \overline{X}$ if and only if there exists a sequence $(x_n)_{n \in \N}$ in $X$ such that $x_n \to x$.
For $n \ge 0$, set $x_n = x_* + \frac{n}{n+1}(x - x_*)$.
Then $d(x_*, x_n) = N(\frac{n}{n+1}(x - x_*)) = \frac{n}{n+1}r < r$.
So $(x_n)_{n \in \N} \in B(x_*, r)^{\N}$ and $x_n \to x$. This implies $x \in \overline{B(x_*, r)}$.
\end{proof}

\subsection{Dense Subsets}
\begin{definition}{Dense Subset}{dense_def}
Let $X, Y \subset E$ be subsets.
\begin{itemize}
    \item $X$ is dense in $E$ if $\overline{X} = E$.
    \item $X$ is dense in $Y$ if $Y \subset \overline{X}$.
\end{itemize}
\end{definition}
\begin{example}{Dense Sets}{dense_ex}
\begin{itemize}
    \item $\overline{\Z} = \Z$, so $\Z$ is not dense in $\R$.
    \item $\overline{\Q} = \R$, so $\Q$ is dense in $\R$.
    \item Let $E = C^0([0,1], \R)$ with the norm $||\cdot||_\infty$. Let $P$ be the subspace of polynomials. By the Weierstrass approximation theorem, $\overline{P} = E$.
\end{itemize}
\end{example}

\section{Lecture 3: Continuous Maps}
Let $(E, N_E)$ and $(F, N_F)$ be $\mathbb{K}$-NVS. Let $\Omega \subset E$ be a non-empty open subset. Let $\varphi: \Omega \to F$.

\begin{definition}{Limit of a Map}{limit_map_def}
Let $l \in F$. $l$ is the limit of $\varphi(w)$ when $w \to w_0$ in $\Omega$ (with $w \neq w_0$) if for all $\eps > 0$, there exists $\delta > 0$ such that for all $w \in \Omega$ with $0 < N_E(w - w_0) < \delta$, we have $\varphi(w) \in B_F(l, \eps)$.
Notation: $\lim_{w \to w_0} \varphi(w) = l$ or $\varphi(w) \to l$ as $w \to w_0$.
\end{definition}

\begin{proposition}{Sequential Characterization of Limits}{seq_char_limits}
We have $\lim_{w \to w_0} \varphi(w) = l$ if and only if for every sequence $(w_n)_{n \in \N}$ in $\Omega$ such that $w_n \to w_0$, we have $\varphi(w_n) \to l$.
\end{proposition}

\begin{definition}{Continuity}{cont_def}
\begin{itemize}
    \item $\varphi$ is continuous ($C^0$) at $w_0 \in \Omega$ if $\lim_{w \to w_0} \varphi(w) = \varphi(w_0)$.
    \item $\varphi$ is continuous on $\Omega$ (notation: $\varphi \in C^0(\Omega, F)$) if $\varphi$ is continuous at every $w \in \Omega$.
\end{itemize}
\end{definition}

\begin{example}{}{cont_ex}
Let $f: \R^2 \to \R$ be defined by $f(x,y) = 0$ if $(x,y) = (0,0)$, and $f(x,y) = \frac{xy^2}{x^2+y^4}$ otherwise. It can be proven (as an exercise) that $f \notin C^0(\R^2, \R)$.
\end{example}

\begin{remark}{}{norm_dep}
Continuity of a map depends on the chosen norm.
\end{remark}
\begin{example}{Norm Dependence}{norm_dep_ex}
Let $E = C^0([0,1], \R) = F$.
Let $N_E = ||\cdot||_\infty$ (sup norm) and $N_F = ||\cdot||_{L^1}$ (where $||f||_{L^1} = \int_0^1 |f(x)|dx$).
Let $\Omega = E$, and let $\varphi: E \to F$ be the identity map, $f \mapsto f$.
The map $\varphi: (E, ||\cdot||_\infty) \to (F, ||\cdot||_{L^1})$ is continuous.
However, $\varphi^{-1}: (F, ||\cdot||_{L^1}) \to (E, ||\cdot||_\infty)$ is not continuous.
To see this, consider $f_n(x) = x^n$ on $[0,1]$.
$||f_n||_{L^1} = \frac{1}{n+1} \to 0$, so $f_n \to 0_F$ in $(F, ||\cdot||_{L^1})$.
But $||f_n||_\infty = 1$ for all $n$, so $f_n$ does not converge to $0_E$ in $(E, ||\cdot||_\infty)$.
\end{example}

\begin{proposition}{Operations on Continuous Functions}{cont_ops}
Let $f, g: \Omega \to F$ and $\lambda \in \mathbb{K}$.
\begin{itemize}
    \item If $f$ and $g$ are continuous at $w_0$ (or on $\Omega$), then $f + \lambda g$ is continuous at $w_0$ (or on $\Omega$). Thus, $C^0(\Omega, F)$ is a $\mathbb{K}$-vector subspace of $F^\Omega$.
    \item If $\mathbb{K} = F$ and $f, g$ are continuous at $w_0$ (or on $\Omega$), then $f \cdot g$ is continuous at $w_0$ (or on $\Omega$).
    \item If $\mathbb{K} = F$, $f$ is continuous at $w_0$, and $f(w_0) \neq 0$, then $\frac{1}{f}$ is defined on some $B(w_0, \eps)$ and is continuous at $w_0$. (If $f \neq 0$ on $\Omega$, then $\frac{1}{f}$ is continuous on $\Omega$).
\end{itemize}
\end{proposition}

\begin{proposition}{Topological Characterization of Continuity}{top_char_cont}
Let $E, F$ be NVS and $\Omega \subset E$ open. Let $\varphi: \Omega \to F$. The following assertions are equivalent:
\begin{enumerate}
    \item $\varphi \in C^0(\Omega, F)$.
    \item For every open subset $U \subset F$, $\varphi^{-1}(U)$ is open (in $E$).
\end{enumerate}
(Note: this also applies when $\Omega = E$).
\end{proposition}
\begin{proof}
$(1 \Rightarrow 2)$ Let $U \subset F$ be an open subset. Let $w \in \varphi^{-1}(U)$, so $\varphi(w) \in U$.
Since $U$ is open, there exists $\eps > 0$ such that $B_F(\varphi(w), \eps) \subset U$.
Because $\varphi$ is continuous at $w$, there exists $\delta > 0$ such that $\varphi(B_E(w, \delta) \cap \Omega) \subset B_F(\varphi(w), \eps)$.
Since $\Omega$ is open, there exists $\delta' \le \delta$ such that $B_E(w, \delta') \subset \Omega$.
Then $\varphi(B_E(w, \delta')) \subset B_F(\varphi(w), \eps) \subset U$.
Thus, $B_E(w, \delta') \subset \varphi^{-1}(U)$, which means $\varphi^{-1}(U)$ is open.

$(2 \Rightarrow 1)$ Let $w \in \Omega$ and $\eps > 0$. Then $U = B_F(\varphi(w), \eps)$ is open in $F$.
By assumption, $\varphi^{-1}(U) = \varphi^{-1}(B_F(\varphi(w), \eps))$ is open in $E$.
Since $w \in \varphi^{-1}(U)$, there exists $\delta > 0$ such that $B_E(w, \delta) \subset \varphi^{-1}(U)$.
This means for all $y \in B_E(w, \delta)$, $\varphi(y) \in B_F(\varphi(w), \eps)$, i.e., $d_F(\varphi(w), \varphi(y)) < \eps$. Thus, $\varphi$ is continuous at $w$.
\end{proof}

\subsection{Continuity of Linear Maps}
\begin{definition}{Continuous Linear Maps}{cont_lin_maps_def}
Let $E, F$ be $\mathbb{K}$-NVS.
Let $L(E,F) = \{L: E \to F \mid L \text{ is } \mathbb{K}\text{-linear}\}$. This is a $\mathbb{K}$-vector space.
Let $\mathcal{L}(E,F) = \{L \in L(E,F) \mid L \text{ is continuous on } E\}$. This is a $\mathbb{K}$-vector subspace of $L(E,F)$.
\end{definition}

\begin{theorem}{Characterization of Continuous Linear Maps}{char_cont_lin_maps}
For $L \in L(E,F)$, the following assertions are equivalent:
\begin{enumerate}
    \item $L$ is continuous on $E$.
    \item $L$ is continuous at $0_E$.
    \item There exists $C \ge 0$ such that for all $x \in E$ with $N_E(x) = 1$, we have $N_F(L(x)) \le C$.
    \item There exists $C \ge 0$ such that for all $x \in E$, $N_F(L(x)) \le C N_E(x)$.
    \item $\sup_{x \neq 0} \frac{N_F(L(x))}{N_E(x)} < +\infty$.
\end{enumerate}
\end{theorem}
\begin{proof}
$1 \Rightarrow 2$: Obvious.

$2 \Rightarrow 3$: For $\eps = 1$, since $L$ is continuous at $0_E$ and $L(0_E) = 0_F$, there exists $\delta > 0$ such that for all $x \in B_E(0_E, \delta)$, we have $L(x) \in B_F(0_F, 1)$.
Let $x \in E$ with $N_E(x) = 1$. Let $\tilde{x} = \frac{\delta}{2}x$. Then $N_E(\tilde{x}) = \frac{\delta}{2} < \delta$, so $\tilde{x} \in B_E(0_E, \delta)$.
Thus, $L(\tilde{x}) \in B_F(0_F, 1)$, meaning $N_F(L(\frac{\delta}{2}x)) \le 1$.
By linearity, $\frac{\delta}{2}N_F(L(x)) \le 1$, which implies $N_F(L(x)) \le \frac{2}{\delta}$. We set $C = \frac{2}{\delta}$.

$3 \Rightarrow 4$: Let $x \in E \setminus \{0\}$. Then $\tilde{x} = \frac{x}{N_E(x)}$ has norm equal to $1$.
So $N_F(L(\tilde{x})) \le C$.
By linearity, $N_F\left(\frac{L(x)}{N_E(x)}\right) \le C$, which implies $N_F(L(x)) \le C N_E(x)$.
(For $x = 0_E$, $N_F(L(0_E)) = 0 = C \cdot 0$).

$4 \Leftrightarrow 5$: By (4), the set $\{\frac{N_F(L(x))}{N_E(x)} \mid x \neq 0\}$ is bounded above by $C$, so its supremum is finite. Let this supremum be denoted as $\Vert L \Vert$.
Then $\Vert L \Vert = \inf\{C \ge 0 \mid \forall x \in E, N_F(L(x)) \le C N_E(x)\}$.

\textbf{($5 \Rightarrow 1$):} Assume $\sup_{x \neq 0} \frac{N_F(L(x))}{N_E(x)} = |||L||| < +\infty$. Then for all $x \in E$, $N_F(L(x)) \le |||L||| N_E(x)$. For any $x_0 \in E$ and $x \in E$, by linearity we have $N_F(L(x) - L(x_0)) = N_F(L(x - x_0)) \le |||L||| N_E(x - x_0)$. As $x \to x_0$, $N_E(x - x_0) \to 0$, which implies $N_F(L(x) - L(x_0)) \to 0$. Thus, $L$ is continuous at any $x_0 \in E$, meaning $L$ is continuous on $E$. This completes the cycle of equivalences.
\end{proof}

\begin{proposition}{Characterization of the Triple Norm}{triple_norm_char}
Let $\mathcal{N}_L = \left\{ \frac{N_F(L(x))}{N_E(x)} \mid x \in E, x \neq 0 \right\} \subset [0, +\infty)$.
Then $\mathcal{N}_L = \{ N_F(L(\tilde{x})) \mid \tilde{x} \in E, N_E(\tilde{x}) = 1 \}$.
Let $\mathcal{C} = \{ C \ge 0 \mid \forall x \in E, N_F(L(x)) \le C N_E(x) \}$ be the set of upper bounds.
Then $\sup \mathcal{N}_L = \inf \mathcal{C}$.
\end{proposition}

\subsection{Uniform Continuity and Lipschitz Continuity}
Let $\varphi: \Omega \to F$.
\begin{definition}{Uniform Continuity}{unif_cont_def}
$\varphi$ is \textbf{uniformly continuous} if:
$\forall \eps > 0, \exists \delta > 0$ such that $\forall x, x' \in \Omega$, if $d_E(x, x') < \delta$ then $d_F(\varphi(x), \varphi(x')) < \eps$.
\end{definition}

\begin{definition}{Lipschitz Continuity}{lipschitz_def}
$\varphi$ is \textbf{Lipschitz} of parameter $\alpha \ge 0$ (or $\alpha$-Lipschitz) if:
$\forall x, x' \in \Omega, d_F(\varphi(x), \varphi(x')) \le \alpha d_E(x, x')$.
\end{definition}

\begin{example}{Lipschitz Maps}{lipschitz_ex}
\begin{enumerate}
    \item By the reverse triangular inequality, $\forall x, x' \in E, |N(x) - N(x')| \le N(x - x')$. Thus, $N: E \to \R$ is 1-Lipschitz.
    \item Let $L \in \mathcal{L}(E, F)$. Then $|||L||| = \sup_{x \neq 0} \frac{N_F(L(x))}{N_E(x)} < +\infty$.
    For all $x, x' \in E$:
    $d_F(L(x), L(x')) = N_F(L(x) - L(x')) = N_F(L(x - x')) \le |||L||| N_E(x - x')$.
    Thus, $L$ is $|||L|||$-Lipschitz.
\end{enumerate}
\end{example}

\begin{proposition}{Implications of Continuity Types}{cont_implications}
For $\varphi: \Omega \to F$:
$\varphi \text{ is Lipschitz} \Rightarrow \varphi \text{ is uniformly continuous (on } \Omega) \Rightarrow \varphi \in C^0(\Omega, F)$.
\end{proposition}
\begin{remark}{}{}
These implications are strict.
\begin{itemize}
    \item $f: [0, +\infty) \to \R$, $x \mapsto \sqrt{x}$ is uniformly continuous on $[0, +\infty)$ but not Lipschitz.
    \item $f: \R \to \R$, $x \mapsto x^2$ is continuous on $\R$ but not uniformly continuous on $\R$.
\end{itemize}
\end{remark}

\subsection{Properties of the Triple Norm}
\begin{definition}{Triple Norm}{triple_norm_def}
For $L \in \mathcal{L}(E, F)$, the triple norm (or operator norm) of $L$ induced by $N_E$ and $N_F$ is:
$|||L||| = \sup_{x \neq 0} \frac{N_F(L(x))}{N_E(x)} \in [0, +\infty)$.
\end{definition}

\begin{proposition}{Properties of Triple Norm}{triple_norm_prop}
\begin{enumerate}
    \item For $L \in \mathcal{L}(E, F)$ and $\forall x \in E$, $N_F(L(x)) \le |||L||| N_E(x)$.
    \item $|||\cdot|||$ is a norm on $\mathcal{L}(E, F)$.
\end{enumerate}
\end{proposition}
\begin{proof}
Let us prove that $|||\cdot|||$ is a norm on $\mathcal{L}(E, F)$.
\begin{itemize}
    \item \textbf{Positive definiteness:} Assume $|||L||| = 0$. Then for all $x \in E$, $N_F(L(x)) = 0$, which implies $L(x) = 0_F$. Thus $L = 0_{\mathcal{L}(E, F)}$.
    \item \textbf{Absolute homogeneity:} Let $L \in \mathcal{L}(E, F)$ and $\lambda \in \mathbb{K}$.
    $|||\lambda L||| = \sup_{x \neq 0} \frac{N_F((\lambda L)(x))}{N_E(x)} = \sup_{x \neq 0} \frac{|\lambda| N_F(L(x))}{N_E(x)} = |\lambda| \sup_{x \neq 0} \frac{N_F(L(x))}{N_E(x)} = |\lambda| |||L|||$.
    \item \textbf{Triangular inequality:} Let $x \in E$ and $L, \tilde{L} \in \mathcal{L}(E, F)$.
    $N_F((L + \tilde{L})(x)) = N_F(L(x) + \tilde{L}(x)) \le N_F(L(x)) + N_F(\tilde{L}(x)) \le |||L||| N_E(x) + |||\tilde{L}||| N_E(x) = (|||L||| + |||\tilde{L}|||) N_E(x)$.
    Thus, $|||L||| + |||\tilde{L}|||$ is an upper bound for $\frac{N_F((L + \tilde{L})(x))}{N_E(x)}$, which implies $|||L + \tilde{L}||| \le |||L||| + |||\tilde{L}|||$.
\end{itemize}
\end{proof}

\begin{proposition}{Unbounded Linear Maps}{unbounded_lin_map}
$L \in L(E,F)$ is not bounded on $B_E(0,1)$ if and only if $L \notin \mathcal{L}(E, F)$ if and only if there exists a sequence $(x_n)_{n \in \N}$ in $E$ such that $N_E(x_n) \to 1$ and $N_F(L(x_n)) \to +\infty$.
\end{proposition}
\begin{proof}
Left as an exercise.
\end{proof}

\begin{proposition}{Computing the Operator Norm}{op_norm_compute}
For $L \in \mathcal{L}(E, F)$ and $C \ge 0$:
\begin{enumerate}
    \item If for all $x \in E$, $N_F(L(x)) \le C N_E(x)$, then $|||L||| \le C$.
    \item If there exists a sequence $(x_n)_{n \in \N}$ in $E$ such that $N_E(x_n) \to 1$ and $N_F(L(x_n)) \to C$, then $|||L||| \ge C$.
\end{enumerate}
(If both conditions are satisfied, then $|||L||| = C$).
\end{proposition}

\begin{example}{Operator on $l^1$}{op_l1_ex}
Let $l^1(\N) = \{ \underline{x} = (x_n)_{n \in \N} \in \R^\N \mid \Vert \underline{x} \Vert_{l^1} < +\infty \}$, where $\Vert \underline{x} \Vert_{l^1} = \sum_{n \ge 0} |x_n|$.
Let $\lambda = (\lambda_n)_{n \in \N} \in l^\infty(\N)$, so $\Vert \lambda \Vert_\infty = \sup_n |\lambda_n| < +\infty$.
Define $\Lambda: l^1(\N) \to l^1(\N)$ by $\Lambda((x_n)_{n \in \N}) = (\lambda_n x_n)_{n \in \N}$.
We have:
$\Vert \Lambda(\underline{x}) \Vert_{l^1} = \sum_{n \ge 0} |\lambda_n x_n| \le \Vert \lambda \Vert_\infty \sum_{n \ge 0} |x_n| = \Vert \lambda \Vert_\infty \Vert \underline{x} \Vert_{l^1}$.
This implies $\Lambda(\underline{x}) \in l^1(\N)$, so $\Lambda$ is well-defined.
$\Lambda$ is clearly $\R$-linear, so $\Lambda \in L(l^1(\N), l^1(\N))$.
From the inequality above, we get that $\Lambda \in \mathcal{L}(l^1(\N), l^1(\N))$ and $|||\Lambda||| \le \Vert \lambda \Vert_\infty$.
(Exercise: prove that $|||\Lambda||| = \Vert \lambda \Vert_\infty$ using the second fact from the proposition above).
\end{example}

\section{Lecture 4: Compactness}
Let $(E, N)$ be a NVS over $\mathbb{K}$.

\begin{definition}{Compactness}{compact_def}
A subset $K \subset E$ is \textbf{compact} if $K \neq \emptyset$ and for every sequence $(x_n)_{n \in \N}$ in $K$, there exists a strictly increasing extraction $\varphi: \N \to \N$ such that the subsequence $(x_{\varphi(n)})_{n \in \N}$ converges in $K$.
\end{definition}

\begin{example}{Compact Sets}{compact_ex}
\begin{itemize}
    \item In $E = \R$, let $I = [a, b]$ with $-\infty < a \le b < +\infty$. Then $I$ is compact.
    \item By the Bolzano-Weierstrass theorem, any bounded sequence in $\R$ admits a convergent subsequence. The proof involves defining $S_n = \inf_{k \ge n} x_k$.
    \item The rectangle $R = [a, b] \times [a', b']$ is compact in $\R^2$ (for any norm, e.g., $\Vert \cdot \Vert_\infty$). (Exercise).
\end{itemize}
\end{example}

\begin{proposition}{Compact Implies Closed and Bounded}{compact_closed_bounded}
For $K \subset E$, if $K$ is compact, then $K$ is bounded and closed.
\end{proposition}
\begin{proof}
We prove this by contrapositive: if $K$ is not bounded or $K$ is not closed, then $K$ is not compact.

Assume $K$ is not bounded.
Then for all $n \in \N$, $K \not\subset B(0, n)$. This means there exists $x_n \in K \setminus B(0, n)$, so $N_E(x_n) \ge n$.
Because any convergent sequence in $E$ is bounded, the sequence $(x_n)_{n \in \N}$ (where $N_E(x_n) \ge n$) can never contain a convergent subsequence. Thus, $K$ is not compact.

Assume $K$ is not closed.
By the sequential criterion for closedness, there exists a sequence $(x_n)_{n \in \N}$ in $K$ such that $x_n \to x_\infty \in E$, but $x_\infty \notin K$.
Then any subsequence $(x_{\varphi(n)})_{n \in \N}$ also converges to $x_\infty$. Since $x_\infty \notin K$, the subsequence does not converge in $K$. Thus, $K$ is not compact.
\end{proof}

\begin{remark}{}{}
\textbf{Warning:} The converse (bounded and closed implies compact) is generally false in infinite-dimensional spaces.
\end{remark}

\begin{example}{Closed and Bounded but Not Compact}{not_compact_ex}
Let $l^\infty(\N) = \{ \underline{x} = (x_n)_{n \in \N} \in \R^\N \mid \Vert \underline{x} \Vert_\infty = \sup_{n \in \N} |x_n| < +\infty \}$.
Let $S = \{ \underline{x} \in l^\infty(\N) \mid N(\underline{x}) = 1 \} = N^{-1}(\{1\})$.
$S$ is clearly bounded.
Since $\{1\}$ is closed in $\R$ and $N$ is 1-Lipschitz (hence continuous), $S = N^{-1}(\{1\})$ is closed in $l^\infty(\N)$.
Thus, $S$ is bounded and closed.
However, consider the sequence $(e_n)_{n \in \N}$ where $e_n = (\delta_{n, k})_{k \in \N} = (0, \dots, 0, 1, 0, \dots)$ (the 1 is at the $n$-th position).
For all $n$, $\Vert e_n \Vert_\infty = 1$, so $e_n \in S$.
For $n \neq m$, $\Vert e_n - e_m \Vert_\infty = 1$.
Therefore, no subsequence of $(e_n)_{n \in \N}$ can be Cauchy, and hence no subsequence can converge. Thus, $S$ is not compact.
\end{example}

\begin{theorem}{Riesz's Theorem}{riesz_thm}
In a NVS $(E, N)$, the closed unit ball $\overline{B_E}(0, 1)$ is compact if and only if $\dim E < +\infty$.
\end{theorem}

\begin{proposition}{Images of Compact Sets}{image_compact}
If $\varphi: \Omega \to F$ is continuous on $\Omega$ and $K \subset \Omega$ is compact, then $\varphi(K)$ is compact.
\end{proposition}
\begin{proof}
Assume $K \neq \emptyset$. We set $\mathcal{K} = \varphi(K)$.
Let $(y_n)_{n \in \N} \in \mathcal{K}^\N$. We know that for all $n \in \N$, there exists $x_n \in K$ such that $y_n = \varphi(x_n)$.
Since $K$ is compact and $(x_n)_{n \in \N} \in K^\N$, there exists a strictly increasing extraction $\psi: \N \to \N$ such that $x_{\psi(n)} \to x_\infty$ for some $x_\infty \in K$.
Because $\varphi$ is continuous at $x_\infty \in K \subset \Omega$, we have that $\varphi(x_{\psi(n)}) \to \varphi(x_\infty)$.
This means $y_{\psi(n)} \to y_\infty \in \mathcal{K}$, so $\mathcal{K}$ is compact.
\end{proof}

\begin{corollary}{Extreme Value Theorem}{extreme_value}
If $K \subset E$ is compact and $f: K \to \R$ is continuous, then $f$ attains its extrema on $K$.
\end{corollary}
\begin{proof}
By the previous proposition, $\mathcal{K} = f(K) \subset \R$ is compact.
Therefore, $\mathcal{K}$ is bounded and closed.
Since $\mathcal{K}$ is bounded from above, $\alpha = \sup \mathcal{K} < +\infty$.
Since $\mathcal{K}$ is closed, $\alpha \in \mathcal{K}$.
Thus, there exists $x^* \in K$ such that $f(x^*) = \alpha$.
For every $x \in K$, we have $f(x) \le \alpha = f(x^*)$, i.e., $f(x^*) = \sup_{x \in K} f(x)$.
The proof for the infimum is similar.
\end{proof}

\section{Lecture 5: Compactness II}
Let $E = \R^d$ (for $d \ge 0$) or $\C^d$. We equip $E$ with the sup norm $\Vert x \Vert_\infty = \max_{i=1, \dots, d} |x^i|$.

\begin{proposition}{Component-wise Convergence}{comp_conv}
Let $(x_n)_{n \in \N}$ be a sequence in $E$, with $x_n = (x_n^1, \dots, x_n^d)$ for all $n \in \N \cup \{\infty\}$.
Then $x_n \to x_\infty$ in $(E, \Vert \cdot \Vert_\infty)$ if and only if for all $i = 1, \dots, d$, $x_n^i \to x_\infty^i$ in $\mathbb{K}$.
\end{proposition}

\begin{proposition}{Heine-Borel Theorem in $\R^d$}{heine_borel}
For $K \subseteq E = \R^d$, $K$ is compact if and only if $K$ is bounded and closed.
\end{proposition}
\begin{proof}
The fact that compact implies bounded and closed was proved in Lecture 4.
Now assume $K \neq \emptyset$, bounded, and closed.
Since $K$ is bounded, there exists $\rho > 0$ such that $K \subseteq B(0, \rho)$ for $\Vert \cdot \Vert_\infty$.
This means $K \subseteq [-\rho, \rho]^d \subset \R^d$.
Because closed subsets of a compact set are compact, it is sufficient to prove that $[-\rho, \rho]^d$ is compact.
We prove this by induction on $d \ge 0$.
\begin{itemize}
    \item $d = 0$: Obvious.
    \item $d = 1$: True by the Bolzano-Weierstrass theorem.
    \item $d \ge 2$: We have a natural linear isomorphism $\R^d \to \R^{d-1} \times \R$, sending $(x^1, \dots, x^d)$ to $((x^1, \dots, x^{d-1}), x^d)$.
    Under this isomorphism, the norm corresponds to $\max(\Vert \cdot \Vert_{\R^{d-1}, \infty}, |\cdot|_{\R})$.
    Then $[-\rho, \rho]^d = [-\rho, \rho]^{d-1} \times [-\rho, \rho]$.
    The first factor is compact by the induction hypothesis, and the second is compact in $(\R, |\cdot|)$.
    By the following lemma, the product of two compact sets is compact.
\end{itemize}
\end{proof}

\begin{lemma}{Product of Metric Spaces}{prod_metric}
Let $(X_1, d_1)$ and $(X_2, d_2)$ be two metric spaces.
Let $X = X_1 \times X_2$, and define $d: X^2 \to \R^+$ by $d((x_1, x_2), (y_1, y_2)) = \max(d_1(x_1, y_1), d_2(x_2, y_2))$.
\begin{enumerate}
    \item $d$ is a metric on $X$.
    \item If $K_i \subset X_i$ are non-empty compact subsets for $i=1,2$, then $K = K_1 \times K_2$ is compact in $X$ for $d$.
\end{enumerate}
\end{lemma}
\begin{proof}
Left as an exercise.
\end{proof}

\subsection{Equivalence of Norms}
Let $\mathbb{K} = \R$ or $\C$, and let $E$ be a NVS over $\mathbb{K}$.

\begin{definition}{Weaker Norm}{weaker_norm_def}
Let $N_1, N_2$ be two norms on $E$. We say $N_2$ is \textbf{weaker} than $N_1$ if there exists $C > 0$ such that $N_2 \le C N_1$, meaning $\forall x \in E, N_2(x) \le C N_1(x)$.
\end{definition}

\begin{example}{}{weaker_norm_ex}
\begin{itemize}
    \item On $\R^2$, $N_1 = \Vert \cdot \Vert_\infty$ and $N_2 = \Vert \cdot \Vert_2$. For $x, y \in \R^2$, $\Vert (x,y) \Vert_2 = \sqrt{x^2 + y^2} \le \sqrt{2 \max(x^2, y^2)} = \sqrt{2} \Vert (x,y) \Vert_\infty$. Thus, $\Vert \cdot \Vert_2$ is weaker than $\Vert \cdot \Vert_\infty$.
    \item Let $E = C^1([0,1], \R)$. Consider the norms $N_1(f) = \Vert f \Vert_\infty + \Vert f' \Vert_\infty$ and $N_2(f) = \Vert f \Vert_\infty$.
    Clearly, $N_2(f) \le 1 \cdot N_1(f)$, so $N_2$ is weaker than $N_1$.
    However, $N_1$ is not weaker than $N_2$. Consider $f_n(x) = x^n$. We have $\Vert f_n \Vert_\infty = 1$, but $N_1(f_n) = 1 + n$. There is no constant $C \ge 0$ such that $1+n \le C \cdot 1$ for all $n$, so $N_1$ is not weaker than $N_2$.
\end{itemize}
\end{example}

\begin{proposition}{Properties of Weaker Norms}{weaker_norm_prop}
Let $N_1, N_2$ be norms on $E$.
\begin{enumerate}
    \item If $N_2$ is weaker than $N_1$, then for any sequence $(x_n)_{n \in \N}$ in $E$:
    \begin{itemize}
        \item If $(x_n)_{n \in \N}$ is bounded for $N_1$, it is bounded for $N_2$.
        \item If $(x_n)_{n \in \N}$ converges for $N_1$, it converges for $N_2$.
    \end{itemize}
    \item If $N_2$ is not weaker than $N_1$, then there exists a sequence $(x_n)_{n \in \N}$ in $E$ that is bounded for $N_1$ but not for $N_2$, and similarly for convergence.
\end{enumerate}
\end{proposition}

\begin{definition}{Equivalent Norms}{equiv_norms_def}
$N_1$ and $N_2$ are \textbf{equivalent} (notation: $N_1 \sim N_2$) if $N_1$ is weaker than $N_2$ and $N_2$ is weaker than $N_1$.
This is equivalent to: $\exists c_1, c_2 > 0$ such that $\forall x \in E$, $c_1 N_1(x) \le N_2(x) \le c_2 N_1(x)$.
\end{definition}

\begin{proposition}{}{equiv_rel}
Equivalence $\sim$ is an equivalence relation on the set of all norms on $E$.
\end{proposition}

\begin{example}{}{equiv_norm_ex}
On $\R^2$, $\Vert \cdot \Vert_\infty \le \Vert \cdot \Vert_2 \le \sqrt{2} \Vert \cdot \Vert_\infty$, so $\Vert \cdot \Vert_\infty \sim \Vert \cdot \Vert_2$.
\end{example}

\begin{theorem}{Equivalence of Norms in Finite Dimensions}{equiv_finite_dim}
On any finite-dimensional vector space $E$, all norms are equivalent.
\end{theorem}

\begin{lemma}{Reduction to Real Scalars}{complex_to_real}
Let $E$ be a $\C$-vector space, and $N$ a norm on $E$.
\begin{enumerate}
    \item Let $E|_\R$ denote $E$ viewed as an $\R$-vector space (restriction of scalars). Then $N$ is a norm on $E|_\R$.
    \item If $N_1, N_2$ are norms on $E$, then $N_1 \sim N_2$ as norms on $E$ if and only if $N_1 \sim N_2$ as norms on $E|_\R$.
\end{enumerate}
\end{lemma}

\begin{proof}[Proof of Theorem \ref{thm:equiv_finite_dim} for $E = \R^d$]
Assume $d > 0$. Let $N$ be a norm on $E$. Let us prove $N \sim \Vert \cdot \Vert_\infty$.
For $x \in E$, we can write $x = \sum_{i=1}^d x^i e_i$ using the canonical basis $(e_i)_{i=1}^d$.
By the triangular inequality, $N(x) = N(\sum_{i=1}^d x^i e_i) \le \sum_{i=1}^d |x^i| N(e_i)$.
Let $C = \sum_{i=1}^d N(e_i) > 0$. Then $N(x) \le \Vert x \Vert_\infty C$, so $N$ is weaker than $\Vert \cdot \Vert_\infty$.

Now we prove the opposite.
Let $S = \{ x \in \R^d \mid \Vert x \Vert_\infty = 1 \}$.
Because $\Vert \cdot \Vert_\infty: (\R^d, \Vert \cdot \Vert_\infty) \to \R$ is 1-Lipschitz and therefore continuous, $S = \Vert \cdot \Vert_\infty^{-1}(\{1\})$ is closed.
Since $S$ is also bounded, $S$ is compact in $(\R^d, \Vert \cdot \Vert_\infty)$.
Since $N$ is weaker than $\Vert \cdot \Vert_\infty$, we have $|N(x) - N(y)| \le N(x-y) \le C \Vert x-y \Vert_\infty$.
Thus, $N: (\R^d, \Vert \cdot \Vert_\infty) \to \R$ is Lipschitz, and therefore continuous.
Since $S$ is compact and $N$ is continuous, $N$ attains its infimum on $S$.
So there exists $x^* \in S$ such that for all $x \in S$, $N(x^*) \le N(x)$.
Since $x^* \in S$, $x^* \neq 0$, which implies $\nu = N(x^*) > 0$.
For any $y \neq 0 \in E$, $\frac{y}{\Vert y \Vert_\infty} \in S$, so $\nu \le N\left(\frac{y}{\Vert y \Vert_\infty}\right)$.
By linearity, $\Vert y \Vert_\infty \le \frac{1}{\nu} N(y)$, meaning $\Vert \cdot \Vert_\infty$ is weaker than $N$.
Thus, $N \sim \Vert \cdot \Vert_\infty$.
\end{proof}

\begin{corollary}{General Finite Dimensional Case}{equiv_finite_gen}
Let $E$ be a finite-dimensional vector space over $\mathbb{K}$ ($\R$ or $\C$). All norms are equivalent on $E$.
\end{corollary}
\begin{proof}
Let $N_1, N_2$ be norms on $E$, and let $d = \dim E > 0$.
Choose a basis $(e_i)_{i=1}^d$ for $E$. The map $\psi: \mathbb{K}^d \to E$ defined by $\psi((x_i)_{i=1}^d) = \sum_{i=1}^d x_i e_i$ is a vector space isomorphism.
Define norms on $\mathbb{K}^d$ by $\tilde{N}_j = N_j \circ \psi$ for $j=1,2$.
By the previous theorem, $\tilde{N}_1 \sim \tilde{N}_2$ on $\mathbb{K}^d$. Because $\psi$ is an isomorphism, this implies $N_1 \sim N_2$ on $E$.
\end{proof}

\begin{corollary}{Topological Properties in Finite Dimensions}{top_finite_dim}
On a finite-dimensional NVS, there is a unique notion of:
\begin{itemize}
    \item open / closed / bounded / compact subsets.
    \item convergence of sequences.
    \item continuity of maps.
\end{itemize}
In finite dimensions, there is no need to specify the norm for these topological properties.
\end{corollary}


\end{document}